% appendix_persistent.tex: the stopping model with a persistent income state: calibration, the run-off, the window ratios (2026-09-10).
% Placed inside Appendix \ref{appendix:mcliquidity}, after the subsection "Calibration to the level of the payment margin".
\subsection{Persistent income: what it changes and what it does not}\label{appendix:persistent}

The model of Tables \ref{tab:mc_runoff_grid} and \ref{tab:mc_runoff_grid2} draws income independently each month. Calibrated to the panel's lifetime charge-off rate, it places $0.93$ of lifetime default in the first year of the loan, against the panel's $0.13$, and its payment elasticity within twelve months is $4.4$ to $5.4$, against the panel's $0.45$. This subsection replaces the independent draw with a persistent income state, calibrates the model to the lifetime rate and to the first-year share, and reports the timing of default, the payment and term elasticities, the run-off of the payment elasticity by loan age, and the term-to-payment ratios $R_{12}$, $R_{24}$ and $R_{36}$ of Table \ref{tab:within_by_group}. Replication: \path{Code/local/mc_persistent_model.R}, \path{Code/exhibits/make_persistent_tables.R}, \path{results/mc_persistent_*.csv}.

\paragraph{The model with an income state.} Log income is the sum of a persistent state and a transitory shock,
\begin{equation}\label{eq:persistent_income}
\log y_t = z_t + \varepsilon_t, \qquad z_{t+1} = \phi\, z_t + \nu_{t+1}, \qquad \nu_{t+1} \sim N(0, s_z^2), \qquad \varepsilon_t \sim N(0, s_\varepsilon^2),
\end{equation}
with $\phi$ the monthly persistence, $s_z$ the standard deviation of the innovation and $s_\varepsilon$ that of the transitory shock, both i.i.d. Median income is one, the payment is a constant $r$ in units of median income, utility is logarithmic, and the borrower starts at her persistent level, $z_1 = 0$, since the lender sets the payment against income at origination; the cross-section of $r$ represents the population's heterogeneity in the payment-to-income ratio. The value of standing $K_t$ of equation \eqref{eq:Krec}, with the balance-scaled penalty of Appendix \ref{appendix:balancepenalty} and a logistic taste shock of scale $s$ on the default margin, becomes a function of the state,
\begin{equation}\label{eq:persistent_K}
K_t(z) = \sigma + \pi_t + \delta s_x \sum_{z'} P(z, z')\, \mathbb{E}_{\varepsilon}\Big[\, s \log\big(1 + e^{(K_{t+1}(z') - g(y'))/s}\big) \,\Big], \qquad K_T(z) = \sigma A_T + \pi_T,
\end{equation}
where $g(y) = u(y) - u(y - r)$ is the utility cost of paying, infinite when $y \le r$, so that default is forced there; $P(z, z')$ is the transition of the discretized state; $\delta = e^{-\rho/12}$ at the annual rate $\rho$; $s_x = e^{-h/12}$ is survival to the next payment at the exit hazard $h = 0.22$ per year; $\pi_t = \kappa_B (b_t - \delta s_x b_{t+1})$ is the one-time balance penalty with $b_t$ the annuity balance at a six percent note rate; and $A_T = \sum_{k=0}^{H-T}\delta^k$ discounts the penalty flow $\sigma$ over the fixed horizon $H = 120$ months. The term inside the expectation is the expected surplus from standing under the taste shock, in place of the positive part in \eqref{eq:Krec}. The borrower defaults at $t$ with probability $p_t(z) = \mathbb{E}_\varepsilon[\Lambda((g(y) - K_t(z))/s)]$, $\Lambda$ the logistic distribution function, and the population over the state evolves as $\mu_{t+1} = s_x P'(\mu_t (1 - p_t))$, so that every moment is a deterministic function of the parameters. The conventions are those of the i.i.d.\ model: $s = 0.5\sigma$, $\log r \sim N(\log 0.1, 0.47^2)$, and $T \in \{48, 60, 72\}$ with weights one quarter, one half, one quarter. The payment elasticity of default within $K$ months is the population mean of $\partial d_K/\partial \log r$ over the mean of $d_K$, the term elasticity is the central difference in $\log T$ at a fixed payment, $R_K$ is their ratio, and the run-off is the elasticity of the bin hazard among survivors on 72-month contracts, 48--60 months over 0--6. At $\phi = s_z = 0$ and $s_\varepsilon = 0.5$ the solver reproduces the hazard paths of \path{Code/local/mc_model_recovery.R} (\path{results/mc_persistent_check.csv}); its moments are the second column of Table \ref{tab:persistent_calibration}.

\paragraph{Calibration.} The targets are the panel's lifetime charge-off rate, $0.048$, and the share of lifetime charge-offs within twelve months, $0.126$ (45,946 of 365,312; \path{results/panel_windows_auto.csv}), at $\rho = 0.09$, $\kappa_B = 0$ and $s_\varepsilon = 0.2$; $(s_z, \sigma)$ are fitted at each $\phi$ on a grid from $0.95$ to $0.998$ (\path{results/mc_persistent_calibration.csv}, rows \texttt{share\_fit\_phi*}). At $\phi \le 0.98$ the share is not reached: raising $s_z$ lowers the first-year share but raises the mass of borrowers whose income falls below the payment at some date, default that no $\sigma$ removes, and that mass exceeds the lifetime target first, so the fit stops at shares of $0.23$, $0.18$ and $0.14$ at $\phi = 0.95$, $0.97$ and $0.98$. At $\phi \ge 0.99$ both targets are met, and $\phi = 0.998$ is the fit whose monthly hazard by loan age is closest to the panel's (\path{results/mc_persistent_hazard.csv}). The calibrated process has $\phi = 0.998$ per month, $0.976$ per year; $s_z = 0.147$, an annual innovation standard deviation of $0.50$ and a standard deviation of $z$ after 72 months of $1.16$; and $\sigma = 0.476$. Table \ref{tab:persistent_calibration} sets it beside the i.i.d.\ model, the i.i.d.\ limit of the same solver, and the data. It matches the first-year share ($0.127$), gives a 24-month share of $0.41$ against the panel's $0.36$, and has a monthly hazard that rises tenfold from the first month to the twelfth and is flat after, at $1.39$ times the twelfth-month level in the sixtieth month, where the panel's rises sixfold and is then flat at $1.30$ times; the i.i.d.\ versions place $0.93$ to $0.96$ of lifetime default in the first year with a hazard that falls from the first month.

\begin{table}[H]
\centering
\caption{The stopping model with i.i.d.\ and with persistent income: parameters, targeted and untargeted moments, and the data}
\label{tab:persistent_calibration}
% replication: Code/local/mc_persistent_model.R, Code/exhibits/make_persistent_tables.R; results/mc_persistent_calibration.csv, results/mc_persistent_iid_limit.csv, results/mc_within_grid_v3.csv
\small
\resizebox{0.95\textwidth}{!}{\begin{tabular}{lrrrr}
\toprule
 & \multicolumn{2}{c}{i.i.d.\ income} & Persistent income & Data \\
\cmidrule(lr){2-3}
 & i.i.d.\ model & i.i.d.\ limit of the & (calibrated) & \\
 & (simulated) & persistent solver & & \\
\midrule
\multicolumn{5}{l}{\emph{Parameters}} \\
Annual rate $\rho$ at calibration & 0.10 & 0.09 & 0.09 &  \\
Persistence $\phi$ (monthly) & 0 & 0 & 0.998 &  \\
Innovation s.d.\ $s_z$ & 0 & 0 & 0.147 &  \\
Transitory s.d.\ $s_\varepsilon$ & 0.5 & 0.50 & 0.20 &  \\
Penalty flow $\sigma$ & 0.251 & 0.236 & 0.476 &  \\
Taste-shock scale $s/\sigma$ & 0.5 & 0.50 & 0.50 &  \\
Balance penalty $\kappa_B$ & 0 & 0 & 0 &  \\
\addlinespace
\multicolumn{5}{l}{\emph{Targeted moments}} \\
Lifetime default & 0.048 & 0.048 & 0.048 & 0.048 \\
Share of lifetime default within 12 months$^{a}$ & 0.93 & 0.96 & 0.13 & 0.13 \\
\addlinespace
\multicolumn{5}{l}{\emph{Untargeted moments}} \\
Share within 24 months & 0.96 & 0.98 & 0.41 & 0.36 \\
Payment elasticity, default within 12 months & 4.53$^{b}$ & 4.38 & 4.59 & 0.45 \\
Payment elasticity, lifetime default &  & 4.33 & 2.22 & 0.35 \\
Term elasticity, default within 12 months$^{c}$ &  & +2.48 & +1.30 & -0.31 \\
$R_{12}$, term over payment elasticity$^{c}$ & +0.464 & +0.565 & +0.284 & -0.70 (0.02) \\
$R_{24}$ & +0.473 & +0.562 & +0.388 & +0.55 (0.01) \\
$R_{36}$ & +0.459 & +0.553 & +0.484 &  \\
Run-off ratio (48--60 over 0--6 months)$^{d}$ &  &  & 0.22 & 0.75 \\
Monthly hazard, month 12 over month 1 &  & 0.01 & 10.00 & 6.05 \\
Monthly hazard, month 60 over month 12 &  & 0.30 & 1.39 & 1.30 \\
\bottomrule
\end{tabular}
}
\vspace{0.5em}
\begin{minipage}{0.95\textwidth}
\noindent \small \textit{Notes:} The i.i.d.\ model: the simulated model of Tables \ref{tab:mc_runoff_grid} and \ref{tab:mc_runoff_grid2} with $\sigma$ calibrated to lifetime default at $\rho = 0.10$ (\path{results/mc_within_grid_v3.csv}, \path{mc_runoff_model_grid_h0.22.csv}). i.i.d.\ limit: the solver of this subsection at $\phi = s_z = 0$, $s_\varepsilon = 0.5$ (\path{results/mc_persistent_iid_limit.csv}). Persistent: the calibrated process (\path{results/mc_persistent_calibration.csv}, row \texttt{calibrated}). Data: the auto-loan panel (\path{results/panel_windows_auto.csv}, \path{identification_within_person_by_group.csv}, \path{runoff_model_rho_co_h0.22.csv}, \path{hazard_by_month.csv}). $^{a}$ Targeted in the persistent version only. $^{b}$ The hazard-bin elasticity in the 0- to 6-month bin of Table \ref{tab:mc_runoff_grid}, the i.i.d.\ model's estimator; the other columns report the elasticity of default within twelve months. $^{c}$ Data: within-borrower estimates of Table \ref{tab:within_by_group}; the pooled fixed-effects term elasticity within twelve months is $+0.01$ ($0.05$); standard errors of the ratios by the delta method. $^{d}$ Data: 61- to 72-month contracts, Table \ref{tab:runoff_shape}. The i.i.d.\ limit's run-off is not reported: with $0.96$ of default in the first year, the bin hazards among survivors after the first year are ratios of vanishing masses. The last two rows are ratios of the monthly hazard by loan age.
\end{minipage}
\end{table}

\paragraph{The payment elasticity.} The persistent version's payment elasticity within twelve months is $4.59$, ten times the panel's $0.45$; over the life of the loan it is $2.22$ against $0.35$; and its term elasticity within twelve months is $+1.30$ against the within-borrower $-0.31$. The elasticity does not move with the income process. Over the 34 feasible points of the $(\phi, s_z)$ scan, $\phi$ from $0.95$ to $0.998$ and $s_z$ from $0.03$ to $0.25$ with $\sigma$ calibrated to the lifetime rate at each point, the elasticity within twelve months lies in $[4.12, 4.60]$ while the first-year share runs from $0.998$ to $0.110$ and the run-off ratio from $71$ to $0.21$; at $s_\varepsilon = 0.1$ and $0.3$ with $s_z$ re-fitted to the share it is $4.50$ and $4.73$ (rows \texttt{scan}, \texttt{s\_eps\_*}). The reason is in \eqref{eq:persistent_K}: the value of standing moves with the persistent state by a multiple of $g$ of order $1/(1 - \delta s_x \phi)$, between ten and forty at these persistences, against a taste shock of scale $0.5\sigma$, so default is a threshold crossing in income at every $\phi$, and the elasticity of a threshold-crossing probability is the ratio of the density to the mass at the threshold, $2.1$ at a mass of five percent for a normal variable in units of its standard deviation: an elasticity of $4$ to $5$ corresponds to a dispersion of log income at the threshold of $0.4$ to $0.5$ log points, and $0.45$ to a dispersion near $4.6$, which no process on the scan approaches. The taste-shock scale is the one lever that changes this (rows \texttt{taste\_*}): at $s/\sigma = 1$, $2$ and $4$ the elasticity is $4.59$, $4.54$ and $3.52$ with the share matched, and at $s/\sigma = 6$ it is $0.71$, with $\sigma = 3.2$, a first-year share of $0.30$ that no $s_z$ brings to the target, $R_{12} = 2.84$ and a run-off of $1.66$. The panel's elasticity is not reached with the share matched in either version.

\paragraph{The run-off.} Table \ref{tab:persistent_runoff} reports the run-off ratio on the $(\rho, \kappa_B)$ grid with $\sigma$ re-calibrated to the lifetime rate at each point, and the profile by loan-age bin at $\rho = 0.10$. In the persistent version the ratio is between $0.20$ and $0.22$ at every rate from $0.02$ to $0.8$ and every $\kappa_B$ from $0$ to $0.2$: it moves by $0.015$ across the rate grid and by $0.002$ across the penalty. In the i.i.d.\ model it is between $0.87$ and $0.99$ at every rate from $0$ to $1.2$ (Table \ref{tab:mc_runoff_grid}); the panel's is $0.75$. The insensitivity to the rate is the same in the two versions; the level and the shape are not. The persistent version's elasticity falls from the first bin, to $0.71$ of its initial level in the 6- to 12-month bin, $0.43$ in the 18- to 24-month bin and $0.21$ in the last, where the panel's stays within $0.06$ of its initial level for three years and then falls to $0.75$, and the i.i.d.\ model's declines to $0.77$.

\begin{table}[!ht]
\centering
\caption{The run-off of the payment elasticity: by discount rate and balance penalty, and by loan-age bin}
\label{tab:persistent_runoff}
% replication: Code/local/mc_persistent_model.R, Code/exhibits/make_persistent_tables.R; results/mc_persistent_grid.csv, results/mc_runoff_model_grid_h0.22.csv, results/runoff_model_rho_co_h0.22.csv
\small
\textit{Panel A: the ratio, 48--60 over 0--6 months, by discount rate}\\[0.4em]
\resizebox{0.95\textwidth}{!}{\begin{tabular}{lrrrrrrrrrrr}
\toprule
Annual rate $\rho$ & 0.00 & 0.02 & 0.05 & 0.10 & 0.20 & 0.30 & 0.35 & 0.50 & 0.55 & 0.80 & 1.20 \\
\midrule
i.i.d.\ income & 0.87 &  & 0.93 & 0.96 & 0.94 &  & 0.93 &  & 0.92 & 0.96 & 0.99 \\
Persistent income, $\kappa_B = 0$ &  & 0.22 & 0.22 & 0.22 & 0.21 & 0.21 &  & 0.21 &  & 0.20 &  \\
Persistent income, $\kappa_B = 0.05$ &  & 0.22 & 0.22 & 0.21 & 0.21 & 0.21 &  & 0.21 &  & 0.20 &  \\
Persistent income, $\kappa_B = 0.2$ &  & 0.22 & 0.22 & 0.22 & 0.21 & 0.21 &  & 0.21 &  & 0.20 &  \\
\midrule
Data, 61--72-month contracts & \multicolumn{11}{l}{0.75} \\
\bottomrule
\end{tabular}
}\\[0.8em]
\textit{Panel B: the profile by loan-age bin at $\rho = 0.10$}\\[0.4em]
\resizebox{0.95\textwidth}{!}{\begin{tabular}{lrrrrrrr}
\toprule
Loan-age bin (months) & 0--6 & 6--12 & 12--18 & 18--24 & 24--36 & 36--48 & 48--60 \\
\midrule
\multicolumn{8}{l}{\emph{Payment elasticity of the bin hazard among survivors}} \\
Data, 61--72-month contracts & 0.65 & 0.63 & 0.63 & 0.64 & 0.61 & 0.54 & 0.49 \\
i.i.d.\ income, $\rho = 0.10$ & 4.53 & 4.43 & 4.11 & 4.09 & 4.45 & 4.34 \\
Persistent income, $\rho = 0.10$, $\kappa_B = 0$ & 5.96 & 4.26 & 3.18 & 2.54 & 1.96 & 1.53 & 1.28 \\
\addlinespace
\multicolumn{8}{l}{\emph{Relative to the 0--6 month bin}} \\
Data & 1.00 & 0.97 & 0.97 & 0.98 & 0.94 & 0.83 & 0.75 \\
i.i.d.\ income & 1.00 & 0.98 & 0.91 & 0.90 & 0.98 & 0.96 \\
Persistent income & 1.00 & 0.71 & 0.53 & 0.43 & 0.33 & 0.26 & 0.21 \\
\bottomrule
\end{tabular}
}
\vspace{0.5em}
\begin{minipage}{0.95\textwidth}
\noindent \small \textit{Notes:} Panel A: the i.i.d.\ model from Table \ref{tab:mc_runoff_grid} (\path{results/mc_runoff_model_grid_h0.22.csv}); the persistent version at the calibrated income process with $\sigma$ re-fitted to lifetime default at each $(\rho, \kappa_B)$ (\path{results/mc_persistent_grid.csv}); the data from Table \ref{tab:runoff_shape}. Panel B: the elasticity of the charge-off hazard to the payment among loans surviving to each bin; the data's row is the 61- to 72-month contracts of Figure \ref{fig:runoff_profiles} (\path{results/runoff_model_rho_co_h0.22.csv}).
\end{minipage}
\end{table}

\paragraph{The window ratios.} Table \ref{tab:persistent_windows} reports $R_{12}$, $R_{24}$ and $R_{36}$ on the $(\rho, \kappa_B)$ surface for both versions, with the panel's within-borrower pair. In the i.i.d.\ model the three ratios are the same at each point for $\kappa_B \le 0.2$, $R_{36} - R_{12}$ between $-0.05$ and $+0.19$, because with $0.93$ of default in the first year the windows differ only in months that contain almost no default; at $\kappa_B = 0.5$, where the penalty spreads default over the first two years and the first-year share is $0.50$ at $\rho = 0.10$, the rise reaches $0.27$ at that rate. In the persistent version the ratio rises with the window at every point of the surface: $R_{36} - R_{12}$ is between $0.10$ and $0.26$, and $R_{24} - R_{12}$ between $0.03$ and $0.13$, largest at low rates and high penalties. All three ratios fall with the rate, $R_{12}$ from $0.37$ at $\rho = 0.02$ to $0.015$ at $\rho = 0.8$ with $\kappa_B = 0$, and the balance penalty lowers each; $R_{12}$ is negative at $(\kappa_B, \rho) = (0.05, 0.8)$, $(0.2, 0.5)$ and $(0.2, 0.8)$, and its minimum on the surface is $-0.099$, where $R_{24} = -0.057$. The panel's pair is $R_{12} = -0.70$ ($0.02$) and $R_{24} = +0.55$ ($0.01$), a rise of $1.25$ from the first window to the second. The smallest $R_{12}$ on the surface is $0.60$ above the panel's, 28 standard errors; the largest $R_{24}$, $0.477$ at $(\kappa_B, \rho) = (0, 0.02)$, is $0.08$ below the panel's, 7 standard errors, with $R_{12} = +0.37$ at that point; and the sign change between the two windows occurs on the surface only at $(0.05, 0.8)$ and $(0.2, 0.5)$, with magnitudes below $0.06$. The persistent version produces the feature the i.i.d.\ version does not, a ratio that rises with the window, at one tenth of the panel's amplitude.

\begin{table}[!ht]
\centering
\caption{The term-to-payment ratios $R_{12}$, $R_{24}$ and $R_{36}$ by discount rate and balance penalty, i.i.d.\ and persistent income, with the data's within-borrower pair}
\label{tab:persistent_windows}
% replication: Code/local/mc_persistent_model.R, Code/exhibits/make_persistent_tables.R; results/mc_persistent_grid.csv, results/mc_within_grid_v3.csv, results/identification_within_person_by_group.csv
\scriptsize
\resizebox{0.9\textwidth}{!}{\begin{tabular}{rrrrrrrrrr}
\toprule
 & & \multicolumn{4}{c}{i.i.d.\ income} & \multicolumn{4}{c}{Persistent income} \\
\cmidrule(lr){3-6} \cmidrule(lr){7-10}
$\kappa_B$ & $\rho$ & $R_{12}$ & $R_{24}$ & $R_{36}$ & $R_{36} - R_{12}$ & $R_{12}$ & $R_{24}$ & $R_{36}$ & $R_{36} - R_{12}$ \\
\midrule
0 & 0.02 & +0.544 & +0.529 & +0.504 & -0.040 & +0.372 & +0.477 & +0.562 & +0.190 \\
 & 0.10 & +0.464 & +0.473 & +0.459 & -0.004 & +0.273 & +0.376 & +0.474 & +0.200 \\
 & 0.30 & +0.259 & +0.291 & +0.299 & +0.040 & +0.119 & +0.202 & +0.309 & +0.190 \\
 & 0.80 & +0.050 & +0.094 & +0.127 & +0.077 & +0.015 & +0.043 & +0.117 & +0.102 \\
\addlinespace
0.05 & 0.02 & +0.532 & +0.513 & +0.485 & -0.047 & +0.352 & +0.465 & +0.556 & +0.204 \\
 & 0.10 & +0.431 & +0.436 & +0.428 & -0.003 & +0.251 & +0.363 & +0.466 & +0.215 \\
 & 0.30 & +0.196 & +0.241 & +0.260 & +0.063 & +0.092 & +0.184 & +0.297 & +0.205 \\
 & 0.80 & -0.065 & -0.009 & +0.046 & +0.111 & -0.020 & +0.016 & +0.097 & +0.116 \\
\addlinespace
0.2 & 0.02 & +0.464 & +0.464 & +0.453 & -0.010 & +0.294 & +0.428 & +0.538 & +0.244 \\
 & 0.10 & +0.295 & +0.359 & +0.370 & +0.074 & +0.188 & +0.321 & +0.444 & +0.256 \\
 & 0.30 & -0.054 & +0.057 & +0.127 & +0.182 & +0.021 & +0.130 & +0.263 & +0.242 \\
 & 0.80 & -0.459 & -0.385 & -0.267 & +0.192 & -0.099 & -0.057 & +0.036 & +0.135 \\
\addlinespace
0.5 & 0.02 & +0.136 & +0.269 & +0.331 & +0.195 &  &  &  &  \\
 & 0.10 & -0.135 & +0.022 & +0.135 & +0.270 &  &  &  &  \\
 & 0.30 & -0.468 & -0.383 & -0.241 & +0.227 &  &  &  &  \\
 & 0.80 & -0.397 & -0.612 & -0.651 & -0.254 &  &  &  &  \\
\midrule
\multicolumn{2}{l}{Data, within borrower} & \multicolumn{8}{l}{$R_{12} = -0.70$ (0.02), $R_{24} = +0.55$ (0.01), $R_{24} - R_{12} = 1.25$} \\
\bottomrule
\end{tabular}
}
\vspace{0.5em}
\begin{minipage}{0.95\textwidth}
\noindent \small \textit{Notes:} $R_K$ is the elasticity of default within $K$ months to the contract length over its elasticity to the payment. i.i.d.: the model with $\sigma$ calibrated to lifetime default at $\rho = 0.10$ for each $\kappa_B$, 150,000 simulated loans, the paper's joint regression with cell effects (\path{results/mc_within_grid_v3.csv}). Persistent: the calibrated income process with $\sigma$ re-fitted at each point, population derivatives (\path{results/mc_persistent_grid.csv}, which also holds $\rho = 0.05$, $0.2$ and $0.5$). Data: the within-borrower estimates of Table \ref{tab:within_by_group}, all loans, standard errors by the delta method.
\end{minipage}
\end{table}

\paragraph{What the persistent version establishes.} First, a persistent income state reproduces the timing of default: the first-year share, the 24-month share within $0.05$, and a hazard that rises over the first year and is flat after. Second, it leaves the payment elasticity at $4.1$ to $4.7$ at every persistence, dispersion and taste-shock scale for which the first-year share is matched, ten times the panel's, because default remains a threshold crossing in income; the i.i.d.\ model has the same elasticity for the same reason. Third, the run-off ratio does not move with the rate in either version, and its level and shape in the persistent version, $0.20$ to $0.22$ with a decline from the first bin, are not the panel's $0.75$ with a flat profile for three years, which the i.i.d.\ model reproduces at a first-year share of $0.93$. Fourth, the persistent version separates the windows, by $0.03$ to $0.13$ between $R_{12}$ and $R_{24}$, with the separation moving with $\rho$ and $\kappa_B$; the panel's separation is $1.25$, with a sign change the surface does not contain at any magnitude above $0.06$. The statements elsewhere in this appendix and in the main text that the run-off, the term margin and the equity margin are insensitive to the rate are therefore established in calibrations whose payment elasticity is ten times the panel's, in both versions of the income process, and they are conditional on that calibration. The model is a device for the mechanisms it contains, the option to default later and the balance-scaled penalty, and not a quantitative match to the bureau's margins. The panel's within-borrower window pattern, $-0.70$ at twelve months and $+0.55$ at 24, measured to a few hundredths, is a fact outside the model class at these calibrations; the two moments the class does not reach, the level of the payment elasticity and the amplitude of the window pattern, are the two that depend on the shape of the default margin rather than on the timing of income, and a version in which the default margin is not a threshold in income is the extension they call for. A second reason the calibration cannot be expected to reach the panel's margins is what it draws from: both income processes are calibrated as earnings processes, whereas the model's income is net of non-discretionary expenses, and \citet{fulford_expense_2024} show that expense shocks are at least as frequent and as large as earnings shocks, with shocks above eighty percent of income several times more common than comparable earnings drops. A process measured on earnings alone understates the dispersion and the tail of the net income the default decision responds to, and no measured earnings process can be plugged in to make the model fit. That is a further reason the paper estimates the rate from sufficient statistics that use neither income nor consumption data: the default responses embed the net shock distribution the household actually faces.
